\origpage{75}{GRAPH MINORS. XI}
\noindent an \textit{open disc} is defined similarly. A \textit{painting} $\Gamma$ in a surface $\Sigma$ is a pair $(U, V)$, where $V \subseteq U \subseteq \Sigma$, $V$ is finite, $U - V$ has only finitely many arc-wise connected components, called \textit{edges}, and for every edge $e$, $\bar{e}$ is a closed disc and $\bar{e} \cap V \subseteq bd(\bar{e})$. The components of $\Sigma - U$ are called \textit{regions}, and further terminology is the same as for drawings. If $\Gamma$ is a painting, then $V(\Gamma)$, $E(\Gamma)$ and the incidence relation between them define a hypergraph, and in this sense we may regard $\Gamma$ as a hypergraph.

A painting or drawing $\Gamma$ in $\Sigma$ is \textit{2-cell} if $\Sigma$ is connected and every region is an open disc (and hence $\Gamma$ is not null). Now we wish to define a ``radial drawing'' of a 2-cell painting or drawing. Roughly, we select a point from each region of $\Gamma$, and make it adjacent to each vertex of $\Gamma$ incident with the region, thereby obtaining a bipartite drawing. However, a region may be ``incident more than once'' with the same vertex, and so we must be a little careful. Let $\Gamma$ be a 2-cell painting or drawing in $\Sigma$, and let $\Gamma'$ be a drawing in $\Sigma$. We say that $\Gamma'$ is a \textit{radial drawing} of $\Gamma$ if it satisfies the following conditions:

\hspace*{12pt}(i) $U(\Gamma) \cap U(\Gamma') = V(\Gamma) \subseteq V(\Gamma')$

\hspace*{12pt}(ii) each region $r$ of $\Gamma$ contains a unique vertex, $v_r$ say, of $\Gamma'$; let
\begin{equation*}
W = \{v_r : r \in R(\Gamma)\}
\end{equation*}
\hspace*{12pt}(iii) $\Gamma'$ is bipartite, and $(V(\Gamma), W)$ is a bipartition of $\Gamma'$

\hspace*{12pt}(iv) if $e, f$ are edges of $\Gamma'$ with the same ends $v, v_r$ say, where $v \in V(\Gamma)$ and $r \in R(\Gamma)$, then $\bar{e} \cup \bar{f}$ does not bound a closed disc in $r \cup \{v\}$

\hspace*{12pt}(v) $\Gamma'$ is maximal subject to (i), ..., (iv).

It is not difficult to show that every 2-cell painting or drawing $\Gamma$ has a radial drawing $\Gamma'$, that $\Gamma'$ is 2-cell, and that every region of $\Gamma'$ includes a unique edge of $\Gamma$, unless $|V(\Gamma)| = 1$ and $E(\Gamma) = \emptyset$. Moreover, a radial drawing is unique up to automorphisms of $\Sigma$ which fix $U(\Gamma)$ pointwise. We \del{omit}\fix{give} the proofs of \fix{the first three of} these assertions, which are elementary but rather tedious\fix{, and we omit the proof of the last one}.

\fix{First let $\Gamma$ be a drawing. At a vertex $x$ of a drawing, the edges at $x$ (a loop counted twice) leave $x$ in a cyclic order. We call these ends of edges the \textit{half-edges} at $x$, and the parts of a small neighbourhood of $x$ between two consecutive half-edges the \textit{corners} at $x$. Each corner is connected and misses the drawing, and so it lies in a region of the drawing. Its two \textit{sides} are the half-edges that bound it. If $x$ is in the closure of a region, then some corner at $x$ lies in this region, since near $x$ the points not in the drawing are in the corners. We use several times that a connected set that misses the boundary of a closed disc lies in its interior or outside it.

For the existence, choose a point in each region of $\Gamma$. These points and the vertices of $\Gamma$, with no edges, form a drawing that satisfies (i)--(iv). We show below that every drawing that satisfies (i)--(iv) has at most $|V(\Gamma)| + 2|E(\Gamma)|$ edges. Then some drawing $\Gamma'$ has the most edges among those that satisfy (i)--(iv), and it satisfies (v). For let $\Gamma''$ satisfy (i)--(iv) and contain $\Gamma'$. By (ii) and (iii), $\Gamma''$ has the same vertices as $\Gamma'$, and it has no more edges, and so $\Gamma'' = \Gamma'$.

For the bound, let $\Gamma'$ satisfy (i)--(iv), and let $H$ be the drawing formed by $\Gamma$ and $\Gamma'$, which is a drawing by (i). By (i), each edge of $\Gamma'$ misses $U(\Gamma)$, and so it lies in a region $r$ of $\Gamma$. By (iii) it joins a vertex of $\Gamma$ to a vertex of $\Gamma'$ not in $V(\Gamma)$, which is in the closure of $r$ but not in $U(\Gamma)$, and so is $v_r$ by (ii). In particular each edge of $\Gamma'$ has just one half-edge at a vertex of $\Gamma$. Let $x \in V(\Gamma)$, suppose that there are at least two half-edges of $H$ at $x$, and suppose that both sides of a corner at $x$ are half-edges of $\Gamma'$. The corner lies in a region $r$ of $\Gamma$, and its sides belong to two different edges $f, g$ of $\Gamma'$ that join $x$ to $v_r$. Choose points of $f$ and of $g$ near $x$, and join them by an arc through the corner. With the rest of $f$ and $g$, this arc forms an $O$-arc in $r$, which bounds a closed disc in $r$, since $r$ is an open disc. The part of the corner cut off by the arc contains $x$, which is not in $r$, and so it lies outside this disc. Together they form a closed disc in $r \cup \{x\}$ bounded by $\bar{f} \cup \bar{g}$, contrary to (iv). So the half-edge that follows a half-edge of $\Gamma'$ in the cyclic order at $x$ is a half-edge of $\Gamma$. Map each edge of $\Gamma'$ to the half-edge that follows its half-edge at its end $x$ in $V(\Gamma)$, or to $x$ itself if there is only one half-edge at $x$. This map is one-to-one, since the map taking each half-edge at $x$ to the next one is one-to-one, and since at most one edge of $\Gamma'$ is at $x$ if there is only one half-edge at $x$. Its values are among the vertices and the $2|E(\Gamma)|$ half-edges of $\Gamma$, which gives the bound.

Now let $\Gamma'$ be a radial drawing of $\Gamma$, and let $H$ be as before. Condition (v) is used in the following way. Let $r$ be a region of $\Gamma$, let $x$ be a vertex of $\Gamma$, and let $K$ be a line in $r \cup \{x\}$ with ends $v_r$ and $x$ that meets $U(\Gamma')$ only in its ends. Then $\Gamma'$ with $K$ added as an edge satisfies (i)--(iii), and so by (v) it fails (iv). That is, $K$ and the closure of some edge $h$ of $\Gamma'$ with ends $x$ and $v_r$ bound a closed disc in $r \cup \{x\}$. The interior of this disc contains no point of $U(H)$. For the interior lies in $r$, and so it misses $U(\Gamma)$, and it contains no vertex of $\Gamma'$, since the only vertices of $\Gamma'$ in $r \cup \{x\}$ are $v_r$ and $x$. Let $g$ be an edge of $\Gamma'$ that meets the interior. It misses the boundary of the disc, since it misses $K$ and $h$, and so it lies in the interior, and its ends are vertices of $\Gamma'$ in the disc, which are $x$ and $v_r$. Then $g$ splits the disc into two closed discs, and one of them is bounded by $\bar{g} \cup \bar{h}$ and lies in $r \cup \{x\}$, contrary to (iv).

Now suppose that $\Gamma$ has an edge. Every vertex $x$ of $\Gamma$ is an end of an edge of $\Gamma$. For otherwise a small closed disc around $x$ meets $U(\Gamma)$ only in $x$. Its boundary bounds a closed disc in the region of $\Gamma$ that contains it, since this region is an open disc. These two closed discs have the same boundary and are different, since only one contains $x$, and so their union is open and closed in $\Sigma$, and it is $\Sigma$, since $\Sigma$ is connected. So $U(\Gamma) = \{x\}$, a contradiction. Also the closure of each region $r$ of $\Gamma$ contains a vertex of $\Gamma$. For $r \ne \Sigma$, since $U(\Gamma) \ne \emptyset$, and so the frontier of $r$ is not empty, since $\Sigma$ is connected. It lies in $U(\Gamma)$, and it contains a vertex, or a point of an edge and then also the ends of this edge. We obtain the following facts, in order.

\hspace*{12pt}(a) \textit{For each region $r$ of $\Gamma$ and each vertex $x$ of $\Gamma$ in the closure of $r$, some edge of $\Gamma'$ joins $x$ to $v_r$.} For $r - U(H)$ is the union of the finitely many regions of $H$ inside $r$, and it is dense in $r$, since $U(H) \cap r$ consists of $v_r$ and edges of $\Gamma'$. So $x$ is in the closure of a region $p$ of $H$ inside $r$. Also $v_r$ is in the closure of every region $q$ of $H$ inside $r$. For $q$ is open and non-empty, and $q \ne r$, since $v_r \in r - q$. Since $r$ is connected, it follows that the frontier of $q$ meets $r$. A point of this frontier in $r$ is in $U(\Gamma') \cap r$, which consists of $v_r$ and edges of $\Gamma'$ that have $v_r$ in their closure, and if it is in such an edge then the whole edge, and so $v_r$, is in the closure of $q$. Hence some corner at $x$ and some corner at $v_r$ lie in $p$. Since $p$ is open and connected, there is a line $K$ as above, which runs from $v_r$ through a corner in $p$, then inside $p$, and reaches $x$ through a corner in $p$. Then $h$ joins $x$ to $v_r$.

\hspace*{12pt}(b) \textit{At each vertex $x$ of $\Gamma$, each corner has one side in $\Gamma$ and the other in $\Gamma'$.} We saw in the proof of the bound that the two sides of a corner are not both in $\Gamma'$, if there are at least two half-edges at $x$. If there is only one, it is a half-edge of $\Gamma$, since $x$ is an end of an edge of $\Gamma$. So suppose that both sides of a corner at $x$ are half-edges of $\Gamma$ (the same one, if it is the only half-edge at $x$). The corner lies in a region $p$ of $H$, inside a region $r$ of $\Gamma$, and as in (a) there is a line $K$ as above, with all points other than its ends in $p$, that reaches $x$ inside this corner. Near $x$, the closed disc bounded by $K$ and $\bar{h}$ lies on one side of $K$, between $K$ and $h$ in the cyclic order. Since its interior misses $U(H)$, no half-edge of $H$ at $x$ lies strictly between $K$ and $h$ on this side. So the half-edge of $h$ at $x$ is next to $K$, and it is one of the two sides of the corner, which are half-edges of $\Gamma$, a contradiction.

\hspace*{12pt}(c) \textit{For each region $r$ of $\Gamma$, different corners at $v_r$ lie in different regions of $H$.} For suppose that two of them lie in one region $p$ of $H$. Since $p$ is open and connected, there is an $O$-arc $J$ with $J - \{v_r\} \subseteq p$ that leaves $v_r$ through one of these corners and returns through the other. Since $J \subseteq r$ and $r$ is an open disc, $J$ bounds a closed disc in $r$. No edge of $\Gamma'$ meets the interior of this disc. For such an edge misses $J$, so it would lie in the disc with both its ends, while each edge of $\Gamma'$ has an end in $V(\Gamma)$, outside $r$. But near $v_r$ the disc lies on one side of $J$, and since the two corners are different, each side of $J$ near $v_r$ contains a half-edge of $\Gamma'$ at $v_r$, a contradiction.

\hspace*{12pt}(d) \textit{For each region $p$ of $H$, one or two edges of $\Gamma'$ are in the closure of $p$, and if two, each of them has different regions of $H$ on its two sides.} Let $r$ be the region of $\Gamma$ that includes $p$. As in (a), $v_r$ is in the closure of $p$, and by (a), $v_r$ is an end of an edge of $\Gamma'$, since the closure of $r$ contains a vertex of $\Gamma$. So some corner at $v_r$ lies in $p$, and by (c) just one does. Each edge of $\Gamma'$ in the closure of $p$ has $v_r$ as an end, and the points of $p$ near it close to $v_r$ lie in a corner beside it, which is then the corner in $p$. So each such edge is a side of this corner. If there are two of them, then the two corners beside one of them at $v_r$ are different, and so by (c) they lie in different regions of $H$.

\hspace*{12pt}(e) \textit{The closure of each region $p$ of $H$ includes exactly one edge of $\Gamma$, and contains no vertex of $\Gamma$ other than the ends of this edge.} Let a half-edge of an edge of $\Gamma$ in the closure of $p$ be given. Some corner beside it lies in $p$, and by (b) the other side of this corner is a half-edge of an edge of $\Gamma'$, which is in the closure of $p$. Each edge of $\Gamma'$ has one half-edge at a vertex of $\Gamma$, and at most two half-edges are next to it there, one on each side. If there are two edges of $\Gamma'$ in the closure of $p$, then by (d) only one side of each of them lies in $p$. So in both cases of (d), at most two half-edges of $\Gamma$ are next to these edges of $\Gamma'$ with the corner between them in $p$. An edge of $\Gamma$ in the closure of $p$ gives two of them, its two half-edges, and so there is at most one such edge. There is at least one, since if $f$ is an edge of $\Gamma'$ in the closure of $p$, then at its end in $V(\Gamma)$ a corner beside $f$ lies in $p$, and by (b) the other side of this corner is a half-edge of $\Gamma$. Finally, if a vertex of $\Gamma$ is in the closure of $p$, then some corner at it lies in $p$, and by (b) one side of this corner is a half-edge of $\Gamma$, whose edge is the edge of $\Gamma$ in the closure of $p$.

\hspace*{12pt}(f) \textit{Each region $s$ of $\Gamma'$ includes exactly one edge $e$ of $\Gamma$, and $s$ is the union of $e$ and the regions of $H$ whose closure includes $e$.} This is the third assertion. Since $V(\Gamma) \subseteq U(\Gamma')$, the set $s$ is the union of the edges of $\Gamma$ in $s$ and the regions of $H$ in $s$. By (e), the closure of each of these regions includes exactly one edge of $\Gamma$, and this edge lies in $s$, since it misses $U(\Gamma')$ and meets the closure of $s$. In particular $s$ includes an edge of $\Gamma$. For each edge $f$ of $\Gamma$ in $s$, let $s_f$ be the union of $f$ and the regions of $H$ whose closure includes $f$. It is open, since each point of $f$ has a neighbourhood that meets $U(H)$ only in $f$, and whose points off $f$ lie in regions of $H$ beside $f$. The sets $s_f$ cover $s$, and they are disjoint by (e). Since $s$ is connected, there is only one of them.

\hspace*{12pt}(g) \textit{$\Gamma'$ is 2-cell, whether or not $\Gamma$ has an edge.} First, for each region $r$ of $\Gamma$, $v_r$ is an end of an edge of $\Gamma'$. If $\Gamma$ has an edge, this follows from (a), since the closure of $r$ contains a vertex of $\Gamma$. Otherwise $V(\Gamma) \ne \emptyset$, since $\Gamma$ is 2-cell and so not null. If no edge of $\Gamma'$ ends at $v_r$, then $U(\Gamma')$ meets $r$ only in $v_r$, and a line $K$ as above, from $v_r$ to a vertex of $\Gamma$ on the frontier of $r$ with its other points in $r$, gives an edge $h$ at $v_r$, a contradiction. Next, each region $p$ of $H$ is an open disc. For $p$ lies in a region $r$ of $\Gamma$, which is an open disc. An $O$-arc $F$ in $p$ bounds a closed disc $\Delta$ in $r$. Each point of $U(\Gamma')$ in $r$ lies in the closure of an edge of $\Gamma'$ with an end outside $r$. This closure is connected, misses $F$, and has a point outside $\Delta$, and so it misses $\Delta$. Hence $\Delta$ misses $U(H)$, and so $\Delta \subseteq p$. Thus $p$ is a connected open subset of the open disc $r$ in which every $O$-arc bounds a closed disc, and so $p$ is an open disc, since a connected open subset of the plane with this property is an open disc. If $\Gamma$ has no edge, then $H = \Gamma'$, and we are done.

Otherwise, each edge $e$ of $\Gamma$ has two different regions of $H$ beside it. For suppose that the two corners beside $e$ at an end $x$ of $e$ lie in one region $p$ of $H$. By (b), their other sides are half-edges of edges $f_1, f_2$ of $\Gamma'$, which are in the closure of $p$. Let $y$ be the other end of $e$ ($y = x$ if $e$ is a loop). As in (e), the half-edge of $e$ at $y$ is next to the half-edge of an edge $f$ of $\Gamma'$ in the closure of $p$, with the corner between them in $p$. Suppose first that $f_1 \ne f_2$. Then by (d) these are the only edges of $\Gamma'$ in the closure of $p$, and each has different regions of $H$ on its two sides. So $f$ is $f_1$ or $f_2$, and since an edge of $\Gamma'$ has just one half-edge at a vertex of $\Gamma$, $y = x$, and $e$ is a loop. But then the half-edge of $f$ at $x$ is next to both half-edges of $e$ at $x$, with the corners between them in $p$, and so $p$ is on both sides of $f$, contrary to (d). Now suppose that $f_1 = f_2$. Then $x$ has just two half-edges, one of $e$ and one of $f_1$, and $p$ is on both sides of $f_1$. By (d), $f_1$ is the only edge of $\Gamma'$ in the closure of $p$, and so $f = f_1$, and as before $y = x$. Then $e$ is a loop at $x$, and $x$ has three half-edges, a contradiction. By (f), it follows that each region of $\Gamma'$ is the union of an edge $e$ of $\Gamma$ and the two regions of $H$ beside $e$, which are disjoint open discs whose closures include $e$, and so it is an open disc, since gluing two disjoint open discs along an edge in the frontier of both gives an open disc. So $\Gamma'$ is 2-cell.

Now let $\Gamma$ be a painting. If the closure $\bar{e}$ of some edge $e$ of $\Gamma$ contains no vertex, then $e$ is the only edge and $\Gamma$ has no vertex. For the sets $bd(\bar{f})$ $(f \in E(\Gamma))$, with $V(\Gamma)$ and one more point on each $bd(\bar{f})$ as vertices, form a drawing whose regions are the regions of $\Gamma$ and the sets $\bar{f} - bd(\bar{f})$ $(f \in E(\Gamma))$, which are open discs. The point set of this drawing is connected. For the frontier of an open disc in a surface is connected, and the frontier of each region is not empty, since $\Sigma$ is connected. So if the point set were the union of two disjoint non-empty closed sets $C_1, C_2$, the frontier of each region would lie in $C_1$ or in $C_2$, and adding to each $C_i$ the closures of the regions whose frontier lies in $C_i$ would split $\Sigma$ into two disjoint non-empty closed sets, which is impossible. But $bd(\bar{e})$ is closed and disjoint from the rest of this point set, which is also closed, and so the rest is empty. In this case $\Sigma$ is a sphere, since it is the union of the closed disc $\bar{e}$ and the region $\Sigma - \bar{e}$ of $\Gamma$, which is an open disc. By (i) and (iii), $\Gamma'$ has no edge, since $V(\Gamma) = \emptyset$, and so it has one vertex. So the only region of $\Gamma'$ is $\Sigma$ minus a point, which is an open disc, and it includes $e$.

Otherwise let $D$ be the drawing formed by the sets $bd(\bar{e})$ $(e \in E(\Gamma))$, with vertex set $V(\Gamma)$. Its regions are the regions of $\Gamma$ and the sets $\bar{e} - bd(\bar{e})$ $(e \in E(\Gamma))$, since these sets are connected, open and disjoint, and their union is $\Sigma - U(D)$. So $D$ is 2-cell. Let $\Gamma'$ satisfy (i)--(iv) for $\Gamma$. For each edge $e$ of $\Gamma$, add to $\Gamma'$ a new vertex in $\bar{e} - bd(\bar{e})$, joined by one edge inside $\bar{e}$ to each vertex of $\Gamma$ in $\bar{e}$. This gives a drawing $G$ that satisfies (i)--(iv) for $D$. For (i)--(iii) hold by the construction, and (iv) holds for the new edges since no two of them have the same ends, and for the other edges since the regions of $\Gamma$ are regions of $D$. Since $G$ has at least as many edges as $\Gamma'$, the bound for $D$ bounds the number of edges of $\Gamma'$, and the existence follows as before.

Now let $\Gamma'$ be a radial drawing of $\Gamma$. Then $G$ is a radial drawing of $D$. For let $G''$ satisfy (i)--(iv) for $D$ and contain $G$. The edges of $G''$ in the regions of $\Gamma$, with the vertices of $\Gamma'$, form a drawing that satisfies (i)--(iv) for $\Gamma$ and contains $\Gamma'$, and so by (v) they are the edges of $\Gamma'$. An edge of $G''$ in $\bar{e} - bd(\bar{e})$ joins the new vertex in it to a vertex $x$ of $\Gamma$ in $\bar{e}$. If two such edges had the same ends, they would form an $O$-arc in $\bar{e}$ meeting $bd(\bar{e})$ only in $x$, which bounds a closed disc in $(\bar{e} - bd(\bar{e})) \cup \{x\}$, contrary to (iv). So $G'' = G$. By the case of drawings, each region of $G$ includes a unique edge of $D$, and each region of the drawing formed by $D$ and $G$ is an open disc. Let $s$ be a region of $\Gamma'$. Each edge of $\Gamma$ that meets $s$ lies in $s$, since it misses $U(\Gamma')$ by (i), and each region of $G$ that meets $s$ lies in $s$, since $U(\Gamma') \subseteq U(G)$. Since $V(\Gamma) \subseteq U(\Gamma')$, the set $s$ is the union of the edges of $\Gamma$ in $s$ and the regions of $G$ in $s$. For each edge $e$ of $\Gamma$ in $s$, the union of $e$ and the regions of $G$ that include an edge of $D$ in $e$ is open, as in (f). These sets cover $s$, and they are disjoint, since each region of $G$ includes just one edge of $D$. Since $s$ is connected, there is only one of them. So $s$ includes a unique edge $e$ of $\Gamma$. Moreover, each edge $a$ of $D$ in $e$ has two regions of the drawing formed by $D$ and $G$ beside it, by (g) for $D$ and $G$. Each of these regions lies inside $\bar{e}$ or misses $\bar{e}$, since it is connected and misses $bd(\bar{e})$. One is inside and one is outside, since near a point of $a$ one side of $a$ is inside $\bar{e}$ and the other is outside. The regions outside, for different edges $a$, are different, by (e) for $D$ and $G$. So $s$ is formed by the open disc $\bar{e} - bd(\bar{e})$ and, for each edge $a$ of $D$ in $e$, the edge $a$ and the region beside it outside $\bar{e}$. Adding these one at a time keeps an open disc, since each time we glue an open disc, disjoint from the set built so far, along the edge $a$. So $\Gamma'$ is 2-cell.}\fixnote{fix:radial-proof-added}{The paper omits these proofs. The argument above is the one used in the formalization, which does not use the uniqueness of radial drawings.}

A curve in a surface $\Sigma$ is a continuous function $\phi : [0, 1] \to \Sigma$. For any curve $\phi$, $|\phi|$ denotes $\{\phi(x) : 0 \le x \le 1\}$. A curve $\phi$ is \textit{constant} if $|\phi|$ is a singleton. Now let $Z \subseteq \Sigma$, and let $\phi_0, \phi_1$ be curves. We say $\phi_0$ is \textit{homotopic to} $\phi_1$ in $Z$ if there is a continuous function $\psi : [0, 1] \times [0, 1] \to Z$ such that
\begin{align*}
\psi(x, 0) &= \phi_0(x) && (0 \le x \le 1) \\
\psi(x, 1) &= \phi_1(x) && (0 \le x \le 1) \\
\psi(0, y) &= \phi_0(0) = \phi_1(0) && (0 \le y \le 1) \\
\psi(1, y) &= \phi_0(1) = \phi_1(1) && (0 \le y \le 1).
\end{align*}
In particular, it is necessary that $\phi_0(0) = \phi_1(0)$ and $\phi_0(1) = \phi_1(1)$. We shall assume without proof certain well-known results about homotopy, such as
